% No preamble. Prefix csg-. Scalar classical background identities only.
\section{Classical background composition modulo gauge}
\label{csg-section}

Fix the primitive odd ratio $\ell$ and write
$M=\ell^k$, $R_1=\ell^{s_1}$ and $R_2=\ell^{s_2}$.
Let $U_n(D)$ be the selected ordinary homogeneous Wilson
minimizer at refinement $n$. The ordinary-minimum theorem
of Notebook15 gives a refinement-uniform small global
$C^{1,\alpha}$ comparison ball, a bound
$\|A_{\mathrm o,n}(D)\|_{C^{1,\alpha}}
\le C\|\log D/i\|_\infty$, and gauge equivalence whenever
equality holds in that comparison class.

Define the pre-root-gauge map and its scalar substitution by
\begin{equation}
 P_{M,R}(D)=M_\ell^k(U_{MR}(D)),\qquad
 (\mathcal T_{M,R}E)(D)=E(P_{M,R}(D)).
 \label{csg-maps}
\end{equation}
Here $R=\ell^s$ and $M_\ell^sP_{M,R}(D)=D$.
All expressions include the natural identification of the
nested lattices. The input $D$ lives on the next-coarse
lattice and $P_{M,R}(D)$ on its $R$-times finer current
lattice. The prefix estimate gives a common classical
radius $b_{\mathrm{cl}}$ and constant $C_0$, independent of
$M,R$ and compatible volume, such that
\begin{equation}
 \left\|\frac1i\log P_{M,R}(e^{iw})\right\|_\infty
                \le\frac{C_0}{R}\|w\|_\infty.
 \label{csg-amplitude}
\end{equation}
The operator $\mathcal T$ is a classical background
substitution, not a fluctuation integral.

\subsection{Two minimum comparisons on the same fine graph}
\label{csg-comparison-section}

For a final datum $Z$ put
\begin{equation}
 U_2=U_{MR_1R_2}(Z),\qquad
 D_1=M_\ell^{k+s_1}(U_2)=P_{MR_1,R_2}(Z),\qquad
 U_1=U_{MR_1}(D_1).
 \label{csg-fields}
\end{equation}
Both fine fields live on the same original oriented graph.
In final-$Z$ units its spacing is
$\eta_2=(MR_1R_2)^{-1}$ and the $D_1$ spacing is $R_2^{-1}$.
The definition of $U_{MR_1}(D_1)$ instead uses $D_1$-unit
coordinates and spacing $\eta_1=(MR_1)^{-1}=R_2\eta_2$.
This is a coordinate relabeling of the same links. In four
dimensions the coupling-free Wilson action is the same
plaquette sum in both conventions.

Exact ordinary composition gives
\begin{align}
 M_\ell^{k+s_1}(U_1)&=D_1=M_\ell^{k+s_1}(U_2),\notag\\
 M_\ell^{k+s_1+s_2}(U_1)&=Z=M_\ell^{k+s_1+s_2}(U_2).
 \label{csg-feasibility}
\end{align}
For fixed finite $R_1,R_2$ both fields lie in the needed
comparison balls after shrinking the real $Z$ radius,
uniformly in $M$ and volume. Explicitly, a fine physical
potential $A_1$ in $D_1$ units becomes
\begin{equation}
 A_1^Z(x)=R_2A_1(R_2x),\qquad
 \|A_1^Z\|_{C^{1,\alpha}}
          \le C R_2^{2+\alpha}\|A_1\|_{C^{1,\alpha}}.
 \label{csg-rescaling}
\end{equation}
The reverse conversion decreases the norm up to a fixed
convention constant, and decreases its supremum by $R_2^{-1}$.
By \eqref{csg-amplitude}, $\|\log D_1/i\|\le C_0\|\log Z/i\|/R_2$.
Thus $U_1$ in final units is bounded by
$CR_2^{1+\alpha}\|\log Z/i\|$, while $U_2$ in $D_1$ units
is bounded by $C\|\log Z/i\|/R_2$. These are finitely many
radius restrictions for fixed ratios.

The minimum at $D_1$, with competitor $U_2$, and the minimum
at $Z$, with competitor $U_1$, yield respectively
\begin{equation}
 S(U_1)\le S(U_2),\qquad S(U_2)\le S(U_1).
 \label{csg-comparisons}
\end{equation}
Equality and the ordinary-minimum equality clause give a
fine gauge $g$ with $U_1=U_2^g$. This gauge is supplied by
the small ordinary-to-corrected transfer, the
constraint-preserving flow into the flat slice, and the
inverse transfer from the same strict slice minimizer.
These controlled gauge transformations can be joined to
the identity inside a smaller common neighborhood.
Equal action alone, without that equality clause, would
not justify this conclusion.

Ordinary gauge covariance gives
\begin{equation}
 M_\ell^k(U_1)=\bigl(M_\ell^k(U_2)\bigr)^{g|_{\mathrm{current}}}.
 \label{csg-gauge-covariance}
\end{equation}
Consequently
\begin{equation}
 P_{M,R_1}(P_{MR_1,R_2}(Z))
       \quad\hbox{is gauge equivalent to}\quad
 P_{M,R_1R_2}(Z).
 \label{csg-orbit-composition}
\end{equation}
The two selected representatives need not coincide. The
connecting gauge need not be unique; its restriction to
$D_1$ sites stabilizes $D_1$.

\subsection{The scalar germ identity and its continuation}
\label{csg-identity-section}

\begin{theorem}[Classical background semigroup modulo gauge]
\label{csg-theorem}
Let $E$ be a scalar total functional holomorphic on a
current-angle product ball and locally gauge invariant
there. On finite compatible regulators,
\begin{equation}
 \mathcal T_{MR_1,R_2}(\mathcal T_{M,R_1}E)
                   =\mathcal T_{M,R_1R_2}E
 \label{csg-semigroup}
\end{equation}
as holomorphic germs at zero. The identity extends to every
connected common product domain of the specified branches.
The output $\mathcal T_{M,R}E$ is locally gauge invariant.
In particular, for any fixed finite $q$,
\begin{equation}
 \mathcal T_{MR^{q-1},R}\cdots
 \mathcal T_{MR,R}\mathcal T_{M,R}E
                         =\mathcal T_{M,R^q}E.
 \label{csg-iteration}
\end{equation}
If $R>C_0$ and $b$ lies in the common classical and original
$E$ domains, both sides of \eqref{csg-iteration} are
analytic on the same radius-$b$ ball and agree there.
\end{theorem}
\begin{proof}
On a sufficiently small real $Z$ ball, the gauge path in
\eqref{csg-gauge-covariance} remains inside the original
domain of $E$. Total gauge invariance proves
\eqref{csg-semigroup}. Both compositions are holomorphic;
uniqueness from the real slice proves the germ identity
and its continuation. No gauge identity for separate
activity coefficients is used.

To see local invariance of the output, extend a small
coarse gauge smoothly to prescribed fine point values.
The transformed minimum at $D$ is a small feasible
competitor at $D^v$. Compare with the minimum at $D^v$,
and perform the inverse-gauge comparison at $D$.
Equality follows; the ordinary-minimum equality clause
identifies the two configurations modulo fine gauge.
Their partial averages are gauge equivalent, so evaluating
$E$ gives local invariance. Analyticity continues this
identity from the real overlap.

The identity can now be iterated, proving
\eqref{csg-iteration}. If each ratio is larger than $C_0$,
\eqref{csg-amplitude} maps the radius-$b$ ball into itself
at every substitution; the direct map has the same
property. Thus both sides are analytic on the connected
common $b$ ball. Their germ equality continues there,
even if the real variational comparison required a
smaller radius depending on the finite ratios or on $q$.
This uses no complex minimum assertion.
\end{proof}

\subsection{Use and limitations of a finite iterate}
\label{csg-scope-section}

The direct representative in \eqref{csg-iteration} obeys
the H\"older prefix estimates
\begin{equation}
 \|a_{M,R^q}\|_\infty\le CbR^{-q},\qquad
 \|\Delta a_{M,R^q}\|_\infty\le CbR^{-2q},\qquad
 \|\Delta^2a_{M,R^q}\|_\infty
                         \le CbR^{-(2+\alpha)q},
 \label{csg-iterated-scales}
\end{equation}
with primitive constants independent of $q$. Its complete
supported families follow by applying the same prefix
theorem at endpoint ratio $R^q$.

This permits keeping the elementary $R$ and grouping $J$
fixed before selecting a finite $q$. A separately proved
bound $C(J)R^{-q\alpha}$ could then have a noncircular
parameter order. Such an activity-norm estimate is not
proved by the scalar identity alone.

The canonical families from successive core/tail
transports need not equal the direct family. Their scalar
sums agree; a new justified localization may represent
that common scalar function. No automatic bound for the
old sequential family, or bounded relocalization on a
chosen quotient norm, follows.

Nor does \eqref{csg-iteration} identify sequential quantum
integrals with one integral at ratio $R^q$. Sequential
steps retain their separate gauge factors, normalizations,
windows and newly generated interactions. Those terms do
not disappear from an iterated quantum recursion. This
is a finite classical background semigroup modulo gauge,
uniform in preceding refinement and compatible volume;
infinite iteration and regulator consistency remain open.
